\chapter*{\makebox[\textwidth][c]{摘\quad 要}}
\addcontentsline{toc}{chapter}{摘要}

医疗资料同时包含自由文本、表格、JSON 记录和 PDF 文档。赛题要求智能体完成数据处理、知识构建和自然语言分析，并使各阶段产物可查询、可校验。MediFlow 以 Nexent 为对话入口，通过 MCP 将智能体的任务级规划与 DataMate 的数据集、算子和任务执行连接起来。

系统沿“数据处理→知识构建→分析查询”组织三项任务。任务一保留输入数据的字段、键名和记录边界，将文本治理结果登记为交付数据集。任务二将清洗记录转为实体、关系和三元组，经离线优先的级联抽取、否定识别、可靠性分级和来源证据校验后写入图谱。任务三由统一分析服务承接 Nexent 对话和 Web 工作台入口，根据医疗语义和数据表生成查询计划，在只读连接、允许查询对象清单和执行时限约束下形成分析结果。

三个智能体均采用“智能体选择任务级工具、服务确定阶段内执行路径”的分工。工具调用负责表达任务意图，服务和算子负责完成确定性处理，来源、运行状态和分析结果通过统一标识关联。

在固定留出集上，离线实体识别使用 2\,500 条 CMeEE 样本，F1 为 37.70\%；离线关系抽取使用 1\,793 条 CMeIE 样本，F1 为 14.31\%。任务三另以实际分析库构建 133 道 NL2SQL 题，固定测试集的初始配置执行准确率为 61.3\%，优化配置验证结果为 95.7\%。报告同时给出抽取质量、查询准确率、在线响应时间和数据规模，使三项任务的处理效果与服务能力具有统一的量化依据。

\vspace{0.7em}
\noindent\textbf{关键词：}
医疗数据治理；MCP；智能体；知识图谱；自然语言数据分析；DataMate；Nexent

\clearpage
